\documentclass{article}
\usepackage{hyperref}
\usepackage[margin=0.5in]{geometry}
\usepackage{amsmath}
\usepackage{graphicx}
\graphicspath{ {./assets/} }

\title{PHYSIOL 578: HW 2.1\footnotetext{U-M GPT was used to refine language of answers derived independently.}}
\author{Jack Yang \\ \href{mailto:yhyunjoo@umich.edu}{yhyunjoo@umich.edu}}
\date{September 29, 2026}

\begin{document}
\maketitle

The Hippo pathway scaffold protein RASSF1A was originally identified as a tumor suppressor in lung and breast cancers. Methylation of the RASSF1A gene has been widely associated with poor patient outcomes. To uncover the signaling pathways perturbed by RASSF1A methylation, an unbiased RNA-seq screening approach was performed in several lung and breast cancer cell lines. Interestingly, the TGF$\beta$ signaling pathway was found to be significantly affected in response to RASSF1A deficiency. Based on these analyses, a hypothesis is formed that RASSF1A may function as a signaling link between the Hippo and TGF$\beta$ signaling pathways.

\section{
    Please describe both canonical and non-canonical TGF$\beta$ signaling pathways (0.5 point)
}

\section{
    Please describe the classical Hippo pathway in a mammalian system (0.5 point)
}

\section{
    If RASSF1A positively promotes the TGF$\beta$ signaling,
    what happens to TGF$\beta$-induced canonical signaling cascade and the expression of its profibrogenic gene targets when the RASSF1A gene is methylated?
    What do you expect to observe regarding the key signaling pathways upon the treatment of TGF$\beta$ (0.3 points).
    Please name three profibrogenic genes that are induced by TGF$\beta$ (0.2 point).
    What are the expression levels of these profibrogenic genes when RASSF1A becomes methylated (0.5 point)?
    (Note: methylation is positively correlated with gene silencing)
}

\section{
    You hypothesize that RASSF1A might enhance YAP's interaction with SMAD signaling in cancer cells.
    Please describe experiments to examine the effects of RASSF1A overexpression and knockdown on YAP phosphorylation (0.5 point)
    and on the YAP-SMAD4 interaction following phosphatic acid stimulation (0.5 point).
}
\end{document}